\chapter*{}
\thispagestyle{empty}
%\cleardoublepage

%\thispagestyle{empty}

\begin{titlepage}
 
 
\setlength{\centeroffset}{-0.5\oddsidemargin}
\addtolength{\centeroffset}{0.5\evensidemargin}
\thispagestyle{empty}

\noindent\hspace*{\centeroffset}\begin{minipage}{\textwidth}

\centering
%\includegraphics[width=0.9\textwidth]{imagenes/logo_ugr.jpg}\\[1.4cm]

%\textsc{ \Large PROYECTO FIN DE CARRERA\\[0.2cm]}
%\textsc{ INGENIERÍA EN INFORMÁTICA}\\[1cm]
% Upper part of the page
% 

\vspace{0.15cm}

%si el proyecto tiene logo poner aquí
\includegraphics[scale=0.25]{imagenes/ugr_logo.jpg} 
 %\vspace{0.95cm}
 
\begin{center}
    Gema Correa Fernández
\end{center}

\vspace{0.5in}
 
% Title
{\LARGE\bfseries Integración de Información Geográfica en la Web Semántica\\
}
\noindent\rule[-1ex]{\textwidth}{0.5pt}\\[3.5ex]
{\Large\bfseries Ontología GEOARES\\[4cm]}
\end{minipage}

\vspace{4.75cm}
\noindent\hspace*{\centeroffset}\begin{minipage}{\textwidth}
\centering

%\textbf{Autor}\\ {Gema Correa Fernández}\\[2.5ex]
%\textbf{Directores}\\
%{José Samos Jiménez}\\[2cm]
%\includegraphics[width=0.15\textwidth]{imagenes/tstc.png}\\[0.1cm]
%\textsc{Departamento de Teoría de la Señal, Telemática y Comunicaciones}\\
%\textsc{---}\\

\begin{flushright}
    Trabajo de Fin de Máster para la ``Integración de Información Geográfica \\ en la Web Semántica'' del Departamento de Lenguajes y Sistemas \\ Informáticos de la Escuela Técnica Superior de Ingenierías \\ Informáticas y de Telecomunicaciones de la UGR \\
    
    \vspace{0.2in}
    
    Director: José Samos Jiménez
\end{flushright}

\vspace{4cm}
Septiembre de 2019
\end{minipage}
%\addtolength{\textwidth}{\centeroffset}
\vspace{\stretch{2}}

 
\end{titlepage}






\cleardoublepage
\thispagestyle{empty}

\begin{center} 
	\large{\textbf{Integración de Información Geográfica en la Web Semántica:\\ SUBTÍTULO}}\\
	\vspace{0.25in}
	Gema Correa Fernández
\end{center}


\section*{Resumen}

% Este trabajo tiene como objetivo el análisis de Información Geográfica mediante el uso de las herramientas QGIS y R, herramientas de software libre de amplia difusión. En particular, está centrado en el estudio del Modelo Digital del Terreno. Para entender el ámbito, se contextualizan y desarrollan los conceptos necesarios del área de los Sistemas de Información Geográfica. El Modelo Digital del Terreno es un modelo de datos que permite la realización de análisis sofisticados; estos análisis se pueden realizar tanto en QGIS como en R. En el trabajo se profundizan y se estudian los análisis considerados más relevantes para este ámbito con ambas herramientas, realizando una comparación y evaluación de ambas. A modo de ejemplo se realiza el análisis para la parte occidental de la Vega de Granada. Los datos utilizados y las funciones desarrolladas mediante paquetes existentes de R, se han estructurado en forma de un nuevo paquete. Por el enfoque que se realiza, este trabajo puede resultar de utilidad para aquellas personas que se adentran en este tema.

Este trabajo tiene como objetivo estudiar las herramientas de la Web Semántica que se pueden utilizar para representar e integrar Información Geográfica. En particular, está centrado en valorar las posibles herramientas que ofrece la Web Semántica y en el desarrollo de una prueba de concepto con Información Geográfica procedente de la provincia de Granada. Para entender el ámbito, se contextualizan y desarrollan los conceptos necesarios del área de los Sistemas de Información Geográfica, además de comprender los conceptos clave de la Web Semántica, que ofrece mecanismos muy útiles para la integración de información y estándares de representación de Información Geográfica mediante diversas herramientas. Por el enfoque que se realiza, este trabajo puede resultar de utilidad para aquellas personas que quieren seguir aprendiendo sobre los Sistemas de Información en un ámbito distinto.

\vspace{0.25in}

{\bf Palabras-clave:} Sistema de Información Geográfica, SIG, Web actual, Web Semántica, OWL, RDF, GEOSPARQL, Shapefile, Protegé, GraphBD.

\cleardoublepage


\thispagestyle{empty}


\begin{center} 
	\large{\textbf{TÍTULO EN INGLÉS:\\ SUBTÍTULO EN INGLÉS}}\\
	\vspace{0.25in}
	Gema Correa Fernández
\end{center}

\section*{Abstract} 



\vspace{0.25in}

{\bf Key-words:} 

\chapter*{}
\thispagestyle{empty}

\noindent\rule[-1ex]{\textwidth}{2pt}\\[4.5ex]

Yo, \textbf{Gema Correa Fernández}, alumna del Máster Universitario de Ingeniería Informática de la \textbf{Escuela Técnica Superior de Ingenierías Informática y de Telecomunicación de la Universidad de Granada}, con DNI 75572158-T, autorizo la ubicación de la siguiente copia de mi Trabajo Fin de Máster en la biblioteca del centro para que pueda ser consultada por las personas que lo deseen.\newline

Asimismo, el código fuente del proyecto y esta documentación pueden consultarse en la dirección \url{https://github.com/Gecofer/TFM} una vez defendido el TFM, para que aquellos que lo deseen puedan conocer el proyecto.

\vspace{5cm}

\noindent Fdo: Gema Correa Fernández

\vspace{2cm}

\begin{flushright}
Granada a \today
\end{flushright}


\chapter*{}
\thispagestyle{empty}

\noindent\rule[-1ex]{\textwidth}{2pt}\\[4.5ex]

D. \textbf{José Samos Jiménez}, Profesor del Departamento de Lenguajes y Sistemas Informáticos de la Universidad de Granada.

\vspace{0.5cm}

\textbf{Informa:}

\vspace{0.5cm}

Que el presente trabajo, titulado \textit{\textbf{Integración de Información Geográfica en la Web Semántica, SUBTÍTULO}},
ha sido realizado bajo su supervisión por \textbf{Gema Correa Fernández}, y autorizamos la defensa de dicho trabajo ante el tribunal que corresponda.

\vspace{0.5cm}

Y para que conste, expiden y firman el presente informe en Granada a \today.

\vspace{1cm}

\textbf{El director:}

\vspace{5cm}

\noindent \textbf{José Samos Jiménez}

\chapter*{Agradecimientos}
\thispagestyle{empty}

       \vspace{1cm}

% A mi apoyo incondicional durante casi 17 años, quién me enseñó a luchar en los momentos más díficiles y a no rendirme bajo ningún concepto.\\

% A toda esa gente que ha estado acompañandome durante este tiempo, familia, amigos y profesores. Gracias a vosotros he crecido tanto personal como profesionalmente.





